% !TEX root = print.tex

\chapter{The problem and the result}
\label{ch:problem}

Let $(\Omega, \mu)$ be a measure space, $\iota$ a finite nonempty index set with
$m = \#\iota$ elements and $\kappa$ a second finite index set. Two families of
square-integrable functions on $\Omega$ are given,
\[
  a = (a_k)_{k \in \iota}, \qquad b = (b_p)_{p \in \kappa},
\]
and we look for $n$ points $x_1, \dots, x_n \in \Omega$ and positive weights
$w_1, \dots, w_n$ such that the weighted point evaluations
$\sum_i w_i |f(x_i)|^2$ stay above $\int |f|^2 \, d\mu$ for every $f$ in the span of the
first family and below the squared coefficient norm for every $f$ in the span of the
second one. The two families enter asymmetrically: the first through its size $m$, the
second only through the \emph{effective dimension} of its Gram matrix.

A statement about the weighted sums is a statement about the two matrices
\[
  \sum_i w_i \, a(x_i) a(x_i)^*, \qquad \sum_i w_i \, b(x_i) b(x_i)^*,
\]
so the whole argument is carried out in the Loewner order: $A \preceq B$ means that
$B - A$ is positive semidefinite. Traces of the matrices appearing below are real, and
the formalisation takes real parts wherever a real number is needed; this is left
silent in the statements.

Each statement names the Lean declaration that formalises it. Every statement and every
proof is formalised: the project compiles without \texttt{sorry}, so each result is
verified by the Lean kernel.

\begin{definition}[Gram matrix]
\label{def:gram}
\lean{Discretization.gram}
\leanok
For a family $a : \Omega \to \mathbb{C}^{\iota}$ of square-integrable functions, the
\emph{Gram matrix} $G(a) \in \mathbb{C}^{\iota \times \iota}$ has entries
\[
  G(a)_{k l} = \int_\Omega a_k(x) \overline{a_l(x)} \, d\mu(x).
\]
\end{definition}

\begin{definition}[Effective dimension]
\label{def:effective-dimension}
\leanok
\uses{def:gram}
Let $J$ be a positive definite matrix and $\Lambda > 0$ with $J \preceq \Lambda \cdot 1$.
The \emph{effective dimension} of $J$ relative to $\Lambda$ is
\[
  M = \frac{\Tr J}{\Lambda}.
\]
\end{definition}

The effective dimension is at most $\#\kappa$, and it is much smaller as soon as the
eigenvalues of $J$ decay; this is the quantity the upper frame bound is stated in, and
it is what makes the theorem applicable when the second family is infinite.

\begin{theorem}[Generalized sparsification]
\label{thm:main}
\lean{Discretization.bss_generalized}
\leanok
\uses{def:gram, def:effective-dimension}
Let $a : \Omega \to \mathbb{C}^{\iota}$ and $b : \Omega \to \mathbb{C}^{\kappa}$ be
square-integrable in every coordinate, let $I = G(a)$ and $J = G(b)$ be positive
definite, and let $\Lambda > 0$ satisfy $J \preceq \Lambda \cdot 1$. Put
$m = \#\iota$ and $M = \Tr J / \Lambda$. Then for every $n \geq m$ there are points
$x_1, \dots, x_n \in \Omega$ and weights $w_1, \dots, w_n > 0$ with
\[
  \Bigl(1 - \sqrt{\tfrac{m-1}{n}}\Bigr)^2 \cdot I
    \; \preceq \; \sum_{i=1}^n w_i \, a(x_i) a(x_i)^*
\]
and
\[
  \sum_{i=1}^n w_i \, b(x_i) b(x_i)^*
    \; \preceq \; \Bigl(1 + \sqrt{\tfrac{M-1}{n}}\Bigr)^2 \Lambda \cdot 1 .
\]
\end{theorem}

\begin{proof}
\leanok
\uses{thm:main-normalized, lem:gram-mulvec}
Replace $a$ by $I^{-1/2} a$. Its Gram matrix is $I^{-1/2} I \, I^{-1/2} = 1$ by
\cref{lem:gram-mulvec}, so \cref{thm:main-normalized} applies to it and produces points
and weights with $(1 - r)^2 \cdot 1$ on the left. Conjugating that inequality by
$I^{1/2}$ preserves the Loewner order and turns it into the bound for $a$ with
$(1 - r)^2 \cdot I$ on the left; the second bound does not involve the first family and
is unchanged.
\end{proof}

This is Theorem 3 of Chkifa, Dolbeault, Krieg and
Ullrich~\cite{CDKU}, and it generalizes the sparsification theorem of Batson,
Spielman and Srivastava~\cite{BSS}, which is the case $a = b$ with $I = J = 1$ and
$M = m$. In eigenvalue form the first bound is the factor
$(1 - \sqrt{(m-1)/n})^2 \lambda_{\min}(I)$ of the paper. No side condition on $m$ or on
$M$ is needed; the potential argument of \cref{ch:construction} does need two, and
\cref{ch:edge} removes them.

Read through quadratic forms, the theorem becomes a statement about norms.

\begin{corollary}[Discretization of the $L_2$-norm]
\label{cor:discretization}
\lean{Discretization.exists_discretization}
\leanok
\uses{thm:main-normalized, lem:integral-quadform}
In the situation of \cref{thm:main} with $I = 1$, the same points and weights satisfy,
for every coefficient vector $c \in \mathbb{C}^{\iota}$,
\[
  \Bigl(1 - \sqrt{\tfrac{m-1}{n}}\Bigr)^2 \int_\Omega
    \Bigl| \sum_{k} \overline{c_k} \, a_k(x) \Bigr|^2 d\mu(x)
  \; \leq \; \sum_{i=1}^n w_i \Bigl| \sum_{k} \overline{c_k} \, a_k(x_i) \Bigr|^2
\]
and, for every $c \in \mathbb{C}^{\kappa}$,
\[
  \sum_{i=1}^n w_i \Bigl| \sum_{p} \overline{c_p} \, b_p(x_i) \Bigr|^2
  \; \leq \; \Bigl(1 + \sqrt{\tfrac{M-1}{n}}\Bigr)^2 \Lambda \sum_{p} |c_p|^2 .
\]
\end{corollary}

\begin{proof}
\leanok
\uses{thm:main-normalized, lem:integral-quadform}
A Loewner inequality between matrices is an inequality between quadratic forms, uniformly
in the coefficient vector. For $f(x) = \sum_k \overline{c_k} a_k(x)$ one has
$\sum_i w_i |f(x_i)|^2 = c^* G c$ with $G$ the accumulated matrix, and
$\int_\Omega |f|^2 d\mu = c^* G(a) c = \|c\|^2$ by \cref{lem:integral-quadform} and
$G(a) = 1$. Reading the second bound of \cref{thm:main-normalized} the same way turns its
right-hand side into $(1 + s)^2 \Lambda \|c\|^2$.
\end{proof}

This is Corollary 4 of~\cite{CDKU}. If the second family is the singular basis of the
embedding of a reproducing kernel Hilbert space $H$ into $L_2(\mu)$, that is, an
orthonormal basis of $H$ which is orthogonal in $L_2$, then the right-hand side is the
squared $H$-norm and $\Lambda$ is the squared norm of the embedding.

\chapter{Matrices in the Loewner order}
\label{ch:loewner}

The construction keeps two matrices positive definite and controls them through two
real numbers. This chapter collects what is needed about the Loewner order, the trace
and rank-one updates. Nothing here mentions sampling; all of it is stated for a general
field $\mathbb{C}$ or $\mathbb{R}$ where the C$^*$-algebra structure is not needed, and
none of it is in Mathlib.

\section{Comparison with multiples of the identity}

\begin{lemma}[Crude bound]
\label{lem:le-trace-smul-one}
\lean{Matrix.PosSemidef.le_trace_smul_one}
\leanok
For every positive semidefinite $A$ we have $A \preceq (\Tr A) \cdot 1$.
\end{lemma}

\begin{proof}
\leanok
\uses{lem:le-smul-one}
The trace is the sum of the eigenvalues of $A$, all of which are nonnegative, so every
single eigenvalue is at most $\Tr A$. A Hermitian matrix whose eigenvalues are all at
most $c$ satisfies $A \preceq c \cdot 1$ by \cref{lem:le-smul-one}.
\end{proof}

\begin{lemma}[Eigenvalue bound]
\label{lem:le-smul-one}
\lean{Matrix.IsHermitian.le_smul_one}
\leanok
Let $A$ be Hermitian and $c$ a real number with $\lambda_i(A) \leq c$ for every
eigenvalue. Then $A \preceq c \cdot 1$.
\end{lemma}

\begin{proof}
\leanok
By the spectral theorem $c \cdot 1 - A$ is unitarily conjugate to the diagonal matrix
with entries $c - \lambda_i(A) \geq 0$, and conjugation by a unitary preserves positive
semidefiniteness.
\end{proof}

\begin{lemma}[The inverse is antitone]
\label{lem:inv-antitone}
\lean{Matrix.PosDef.inv_le_inv_of_le, Matrix.PosDef.smul_one_le_of_inv_le}
\leanok
Let $A$ and $B$ be positive definite with $A \preceq B$. Then
$B^{-1} \preceq A^{-1}$. In particular, if $A^{-1} \preceq c \cdot 1$ for some $c > 0$,
then $c^{-1} \cdot 1 \preceq A$.
\end{lemma}

\begin{proof}
\leanok
Both are read off Mathlib's antitonicity of the inverse in a C$^*$-algebra with a star
ordering, applied to the units of $\mathbb{C}^{n \times n}$ with the
$\ell_2$-operator norm; the inverse of $c \cdot 1$ is $c^{-1} \cdot 1$.
\end{proof}

\begin{lemma}[Conjugation by the square root]
\label{lem:sqrt-conjugation}
\lean{CStarAlgebra.inv_smul_le_of_conj_inv_sqrt_le}
\leanok
Let $A$ be a C$^*$-algebra with a star ordering, let $B \in A$ be positive and
invertible, let $X \in A$ and let $t > 0$. If
$B^{-1/2} X B^{-1/2} \preceq t \cdot 1$, then $t^{-1} \cdot X \preceq B$.
\end{lemma}

\begin{proof}
\leanok
Conjugating the hypothesis by $B^{1/2}$ preserves the order. On the left the two
inverse square roots cancel, on the right $t \cdot 1$ becomes $t \cdot B$; dividing by
$t$ gives the claim.
\end{proof}

The next two lemmas are the reason the whole argument can be run on two real numbers.
The first is the lower, the second the upper one.

\begin{lemma}[A bound on the lower potential bounds the matrix]
\label{lem:inv-trace-smul-one-le}
\lean{Matrix.PosDef.inv_re_trace_smul_one_le}
\leanok
For every positive definite $A$,
\[
  \bigl( \Tr A^{-1} \bigr)^{-1} \cdot 1 \; \preceq \; A .
\]
\end{lemma}

\begin{proof}
\leanok
\uses{lem:le-trace-smul-one, lem:inv-antitone}
\cref{lem:le-trace-smul-one} applied to $A^{-1}$ gives
$A^{-1} \preceq (\Tr A^{-1}) \cdot 1$, and \cref{lem:inv-antitone} turns that around.
\end{proof}

\begin{lemma}[A bound on the upper potential bounds the matrix]
\label{lem:inv-trace-mul-smul-le}
\lean{Matrix.PosDef.inv_re_trace_mul_smul_le}
\leanok
For all positive definite $B$ and $J$,
\[
  \bigl( \Tr (J B^{-1}) \bigr)^{-1} \cdot J \; \preceq \; B .
\]
\end{lemma}

\begin{proof}
\leanok
\uses{lem:le-trace-smul-one, lem:sqrt-conjugation, lem:trace-mul-pos}
Write $S = B^{1/2}$. The matrix $S^{-1} J S^{-1}$ is positive semidefinite and has
trace $\Tr (J B^{-1})$, so \cref{lem:le-trace-smul-one} bounds it by
$\Tr (J B^{-1}) \cdot 1$. That trace is positive by \cref{lem:trace-mul-pos}, so
\cref{lem:sqrt-conjugation} applies and gives the claim.
\end{proof}

\begin{lemma}[The two shifts]
\label{lem:shifts}
\lean{Matrix.PosDef.sub_smul_one, Matrix.PosDef.add_smul}
\leanok
Let $A$ be positive definite and $0 < \delta < (\Tr A^{-1})^{-1}$. Then
$A - \delta \cdot 1$ is positive definite. If $B$ and $J$ are positive definite and
$\zeta > 0$, then $B + \zeta \cdot J$ is positive definite.
\end{lemma}

\begin{proof}
\leanok
\uses{lem:inv-trace-smul-one-le}
For the first, write
$A - \delta \cdot 1 = \bigl( A - (\Tr A^{-1})^{-1} \cdot 1 \bigr)
  + \bigl( (\Tr A^{-1})^{-1} - \delta \bigr) \cdot 1$, a positive semidefinite matrix
by \cref{lem:inv-trace-smul-one-le} plus a positive definite one. The second is the sum
of a positive definite matrix and a positive multiple of one.
\end{proof}

\section{Traces of products}

\begin{lemma}[Sign of the trace of a product]
\label{lem:trace-mul-pos}
\lean{Matrix.PosSemidef.trace_mul_nonneg, Matrix.PosDef.re_trace_mul_pos}
\leanok
For positive semidefinite $P$ and $Q$ we have $\Tr (P Q) \geq 0$. If $P$ is positive
definite, $Q$ is positive semidefinite and $Q \neq 0$, then $\Tr (P Q) > 0$.
\end{lemma}

\begin{proof}
\leanok
Write $Q$ as a sum of rank-one matrices $v v^*$. Then
$\Tr (P v v^*) = v^* P v \geq 0$, with strict inequality for $v \neq 0$ when $P$ is
positive definite, and at least one of the $v$ is nonzero when $Q \neq 0$.
\end{proof}

\begin{lemma}[Cauchy--Schwarz for the trace]
\label{lem:trace-cauchy-schwarz}
\lean{Matrix.PosSemidef.norm_trace_mul_sq_le, Matrix.PosSemidef.re_trace_mul_sq_le}
\leanok
Let $Y$ be positive semidefinite and $Z$ any matrix. Then
\[
  \bigl| \Tr (Z Y) \bigr|^2 \; \leq \; \Tr Y \cdot \Tr (Z Y Z^*) .
\]
\end{lemma}

\begin{proof}
\leanok
A positive semidefinite $Y$ makes $\langle P, Q \rangle = \Tr (Q Y P^*)$ a
semi-inner product on matrices. Cauchy--Schwarz for the pair $1$, $Z$ is the
inequality. No commutation of $Y$ and $Z$ is used.
\end{proof}

\section{Rank-one updates}

\begin{lemma}[Sherman--Morrison]
\label{lem:sherman-morrison}
\lean{Matrix.inv_add_smul_vecMulVec}
\leanok
Let $A$ be Hermitian and invertible, $v$ a vector and $t$ a scalar with
$1 + t \, v^* A^{-1} v \neq 0$. Then $A + t \, v v^*$ is invertible and
\[
  \bigl( A + t \, v v^* \bigr)^{-1}
    = A^{-1} - \frac{t}{1 + t \, v^* A^{-1} v} \,
      \bigl( A^{-1} v \bigr) \bigl( A^{-1} v \bigr)^* .
\]
\end{lemma}

\begin{proof}
\leanok
Multiplying the right-hand side by $A + t \, v v^*$ collapses, by Hermitian symmetry of
$A^{-1}$, to $1$ plus a multiple of $v (A^{-1} v)^*$ whose scalar factor vanishes by
the choice of the coefficient.
\end{proof}

\begin{lemma}[Trace of a rank-one update]
\label{lem:trace-sherman-morrison}
\lean{Matrix.trace_inv_add_smul_vecMulVec}
\leanok
In the situation of \cref{lem:sherman-morrison},
\[
  \Tr \bigl( A + t \, v v^* \bigr)^{-1}
    = \Tr A^{-1} - \frac{t}{1 + t \, v^* A^{-1} v} \, v^* A^{-2} v .
\]
\end{lemma}

\begin{proof}
\leanok
\uses{lem:sherman-morrison}
Take traces in \cref{lem:sherman-morrison} and use
$\Tr \bigl( (A^{-1} v)(A^{-1} v)^* \bigr) = v^* A^{-2} v$.
\end{proof}

\begin{lemma}[Positivity of an update]
\label{lem:update-posdef}
\lean{Matrix.PosDef.add_smul_vecMulVec, Matrix.PosDef.sub_smul_vecMulVec}
\leanok
Let $A$ be positive definite, $v$ a vector and $w \geq 0$. Then $A + w \, v v^*$ is
positive definite. If moreover $w \, v^* A^{-1} v < 1$, then $A - w \, v v^*$ is
positive definite.
\end{lemma}

\begin{proof}
\leanok
\uses{lem:sherman-morrison}
The first is a positive definite matrix plus a positive semidefinite one. For the
second, \cref{lem:sherman-morrison} exhibits the inverse of $A - w \, v v^*$ as a
positive definite matrix plus a positive semidefinite one, and a matrix whose inverse is
positive definite is itself positive definite.
\end{proof}

\begin{lemma}[A rank-one matrix is bounded by its norm]
\label{lem:rank-one-crude}
\lean{Matrix.vecMulVec_le_norm_sq_smul_one}
\leanok
For every vector $v$ we have $v v^* \preceq \|v\|^2 \cdot 1$.
\end{lemma}

\begin{proof}
\leanok
\uses{lem:le-trace-smul-one}
The trace of $v v^*$ is $\|v\|^2$, so this is \cref{lem:le-trace-smul-one}.
\end{proof}

\chapter{Averages of quadratic forms}
\label{ch:averages-of-forms}

The bridge between the analytic side of the problem and the matrix algebra is one
identity: the average of a quadratic form along a family is a trace against its Gram
matrix. Read from right to left it says that a trace is an average, which is how the
construction finds a point.

\begin{lemma}[Average of a quadratic form]
\label{lem:integral-quadform}
\lean{Discretization.integral_quadForm, Discretization.integral_re_quadForm}
\leanok
\uses{def:gram}
Let $a : \Omega \to \mathbb{C}^{\iota}$ be square-integrable in every coordinate. Then
for every matrix $Q \in \mathbb{C}^{\iota \times \iota}$,
\[
  \int_\Omega a(x)^* Q \, a(x) \, d\mu(x) = \Tr \bigl( Q \, G(a) \bigr).
\]
\end{lemma}

\begin{proof}
\leanok
Expand both sides as a double sum over $\iota$ and exchange the finite sum with the
integral. The integrand is integrable by Cauchy--Schwarz for $L_2$-functions.
\end{proof}

\begin{lemma}[Gram matrix of a transformed family]
\label{lem:gram-mulvec}
\lean{Discretization.gram_mulVec}
\leanok
\uses{def:gram}
For every matrix $S \in \mathbb{C}^{\iota \times \iota}$ and every family $a$ that is
square-integrable in each coordinate,
\[
  G(S a) = S \, G(a) \, S^* .
\]
\end{lemma}

\begin{proof}
\leanok
Expand both sides entrywise and exchange the finite sums with the integral.
\end{proof}

\begin{lemma}[One good point suffices]
\label{lem:exists-point}
\lean{Discretization.exists_lt_of_integral_lt}
\leanok
Let $f, g : \Omega \to \mathbb{R}$ be integrable with
$\int_\Omega g \, d\mu < \int_\Omega f \, d\mu$. Then $g(x) < f(x)$ for at least one
$x \in \Omega$.
\end{lemma}

\begin{proof}
\leanok
If $f \leq g$ everywhere then $\int f \leq \int g$.
\end{proof}

The paper states the same step as: the set of admissible points has positive measure.
One point is all the construction uses. When the two averages are equal, which happens
in the doubly degenerate case of \cref{ch:edge}, a strict inequality is no longer
available and the following variant takes over.

\begin{lemma}[One good point when the averages agree]
\label{lem:exists-point-le}
\lean{Discretization.exists_le_of_integral_le}
\leanok
Let $f, g : \Omega \to \mathbb{R}$ be integrable and nonnegative with
$\int_\Omega g \, d\mu \leq \int_\Omega f \, d\mu$ and
$\int_\Omega f \, d\mu > 0$. Then there is an $x \in \Omega$ with
$g(x) \leq f(x)$ and $f(x) > 0$.
\end{lemma}

\begin{proof}
\leanok
If $f \leq g$ everywhere then the two averages agree, so the nonnegative function
$g - f$ has average zero and vanishes almost everywhere. At every point where $f$ and
$g$ agree, the assumption forces $f$ to vanish, so $f$ vanishes almost everywhere,
contradicting $\int f > 0$.
\end{proof}

Both the nonnegativity of $g$ and the conclusion $f(x) > 0$ are needed: for
$f = (1, 0)$ and $g = (2, -1)$ on two points with the counting measure both averages
are $1$, and the only point with $g \leq f$ has $f = 0$.

\chapter{Potentials, verifiers and the barrier lemma}
\label{ch:barrier}

Two matrices are carried through the construction: a small one $A$, which must stay
above a multiple of the identity, and a large one $B$, which must stay below one. Each
is controlled by a single real number.

\begin{definition}[The two potentials]
\label{def:potentials}
\lean{Discretization.lowerPotential, Discretization.upperPotential}
\leanok
For a positive definite $A$ and positive definite $B$, $J$, put
\[
  \Phi(A) = \Tr A^{-1}, \qquad \Psi_J(B) = \Tr \bigl( J B^{-1} \bigr).
\]
\end{definition}

The lower potential is large when $A$ has a small eigenvalue, the upper one when $B$ is
small in the directions where $J$ is large. Before each step $A$ is shrunk by
$\delta \cdot 1$ and $B$ is grown by $\zeta \cdot J$; both shifts move the potentials in
the unfavourable direction by an amount that is computed exactly, and that amount is the
budget one new sampling point is allowed to consume.

\begin{lemma}[Effect of the shifts]
\label{lem:shift-effect}
\lean{Discretization.lowerPotential_sub_eq, Discretization.upperPotential_sub_eq}
\leanok
\uses{def:potentials}
Let $A$, $B$ and $J$ be positive definite, let $\delta > 0$ with
$A - \delta \cdot 1$ positive definite and let $\zeta \geq 0$. Then
\[
  \Phi(A - \delta \cdot 1) - \Phi(A)
    = \delta \, \Tr \bigl( (A - \delta \cdot 1)^{-1} A^{-1} \bigr)
\]
and
\[
  \Psi_J(B) - \Psi_J(B + \zeta \cdot J)
    = \zeta \, \Tr \bigl( J B^{-1} J (B + \zeta \cdot J)^{-1} \bigr).
\]
Both right-hand sides are positive for $\delta, \zeta > 0$.
\end{lemma}

\begin{proof}
\leanok
\uses{lem:trace-mul-pos, lem:shifts}
Both identities follow from the resolvent identity
$X^{-1} - Y^{-1} = X^{-1} (Y - X) Y^{-1}$ with the difference of the two matrices being
$\delta \cdot 1$ and $\zeta \cdot J$ respectively. The signs are \cref{lem:trace-mul-pos}
applied to the products of positive definite matrices on the right.
\end{proof}

Whether a candidate point may be used is decided pointwise, without recomputing any
inverse, by the two verifiers.

\begin{definition}[The two verifiers]
\label{def:verifiers}
\lean{Discretization.lowerVerifier, Discretization.upperVerifier}
\leanok
\uses{def:potentials}
Write $Z = (A - \delta \cdot 1)^{-1}$ and $X = (B + \zeta \cdot J)^{-1}$. For a vector
$v$ put
\[
  L_A^{\delta}(v)
    = \frac{v^* Z^2 v}{\Phi(A - \delta \cdot 1) - \Phi(A)} - v^* Z v,
\]
\[
  U_{B}^{\zeta}(v)
    = \frac{v^* X J X v}{\Psi_J(B) - \Psi_J(B + \zeta \cdot J)} + v^* X v .
\]
\end{definition}

\begin{lemma}[Potential of a rank-one update]
\label{lem:update-potential}
\lean{Discretization.lowerPotential_add_smul_vecMulVec,
      Discretization.upperPotential_sub_smul_vecMulVec}
\leanok
\uses{def:potentials}
Let $N$ be positive definite, $v$ a vector and $w \geq 0$. Then
\[
  \Phi \bigl( N + w \, v v^* \bigr)
    = \Phi(N) - \frac{w}{1 + w \, v^* N^{-1} v} \, v^* N^{-2} v .
\]
If $M$ is positive definite and $w \, v^* M^{-1} v < 1$, then
\[
  \Psi_J \bigl( M - w \, v v^* \bigr)
    = \Psi_J(M) + \frac{w}{1 - w \, v^* M^{-1} v} \, v^* M^{-1} J M^{-1} v .
\]
\end{lemma}

\begin{proof}
\leanok
\uses{lem:sherman-morrison, lem:trace-sherman-morrison}
Insert the Sherman--Morrison formula of \cref{lem:sherman-morrison} and take traces as
in \cref{lem:trace-sherman-morrison}.
\end{proof}

\begin{lemma}[Barrier lemma]
\label{lem:barrier}
\lean{Discretization.lowerPotential_update_le, Discretization.upperPotential_update_le}
\leanok
\uses{def:verifiers}
Let $A$ be positive definite, $0 < \delta < \Phi(A)^{-1}$, and let $w > 0$ satisfy
$1/w \leq L_A^{\delta}(v)$. Then $A - \delta \cdot 1 + w \, v v^*$ is positive definite
and
\[
  \Phi \bigl( A - \delta \cdot 1 + w \, v v^* \bigr) \leq \Phi(A).
\]
Let $B$ and $J$ be positive definite, $\zeta > 0$, and let $w > 0$ satisfy
$U_B^{\zeta}(v) \leq 1/w$. Then $B + \zeta \cdot J - w \, v v^*$ is positive definite
and
\[
  \Psi_J \bigl( B + \zeta \cdot J - w \, v v^* \bigr) \leq \Psi_J(B).
\]
\end{lemma}

\begin{proof}
\leanok
\uses{lem:update-potential, lem:update-posdef, lem:shifts}
Both halves are the closed form of \cref{lem:update-potential} combined with one
inequality between real numbers, and positive definiteness is \cref{lem:update-posdef}.
In the upper half the hypothesis on $w$ first has to be turned into positivity of the
Sherman--Morrison denominator $1 - w \, v^* X v$, which is where positive definiteness of
$J$ enters.
\end{proof}

\chapter{The verifiers pass on average}
\label{ch:averages}

A point may be used if the upper verifier stays below the lower one. Such a point exists
because the inequality holds on average.

\begin{lemma}[Average of the lower verifier]
\label{lem:average-lower}
\lean{Discretization.integral_lowerVerifier_gt}
\leanok
\uses{def:verifiers, def:gram}
Let $a$ be square-integrable in every coordinate with $G(a) = 1$, let $A$ be positive
definite and $0 < \delta < \Phi(A)^{-1}$. Then
\[
  \int_\Omega L_A^{\delta}\bigl( a(x) \bigr) d\mu(x)
    \; > \; \frac{1}{\delta} - \Phi(A).
\]
\end{lemma}

\begin{proof}
\leanok
\uses{lem:integral-quadform, lem:trace-cauchy-schwarz, lem:shift-effect}
With $Y = A^{-1}$ and $Z = (A - \delta \cdot 1)^{-1}$, \cref{lem:integral-quadform}
turns the average into $\Tr (Z^2) / (\delta \Tr (Y Z)) - \Tr Z$. Writing
$t = \Tr (Y Z)$ and $u = \Tr (Y Z^2)$, the claim reduces to $\delta t^2 < u$, which is
what \cref{lem:trace-cauchy-schwarz} gives together with $\delta \Phi(A) < 1$.
\end{proof}

\begin{lemma}[Average of the upper verifier]
\label{lem:average-upper}
\lean{Discretization.integral_upperVerifier_lt}
\leanok
\uses{def:verifiers, def:gram}
Let $b$ be square-integrable in every coordinate with $G(b) = J$ positive definite, let
$B$ be positive definite and $\zeta > 0$. Then
\[
  \int_\Omega U_B^{\zeta}\bigl( b(x) \bigr) d\mu(x)
    \; < \; \frac{1}{\zeta} + \Psi_J(B).
\]
\end{lemma}

\begin{proof}
\leanok
\uses{lem:integral-quadform, lem:inv-antitone, lem:shift-effect, lem:trace-mul-pos}
Write $W = B^{-1}$ and $X = (B + \zeta \cdot J)^{-1}$ and let
$E = \Psi_J(B) - \Psi_J(B + \zeta \cdot J)$ be the gain of \cref{lem:shift-effect}. By
\cref{lem:integral-quadform} the average equals
$\Tr (J X J X) / E + \Psi_J(B + \zeta \cdot J)$. Since $X \preceq W$ by
\cref{lem:inv-antitone}, the numerator is at most $\Tr (J W J X) = E / \zeta$, so the
first summand is at most $1/\zeta$, and the second is strictly smaller than $\Psi_J(B)$.
\end{proof}

\begin{lemma}[An admissible point exists]
\label{lem:admissible-point}
\lean{Discretization.exists_admissible_point}
\leanok
\uses{lem:average-lower, lem:average-upper}
In the situation of \cref{lem:average-lower} and \cref{lem:average-upper}, assume
\[
  \frac{1}{\zeta} + \Psi_J(B) \; \leq \; \frac{1}{\delta} - \Phi(A).
\]
Then there is an $x \in \Omega$ with
$U_B^{\zeta}\bigl( b(x) \bigr) < L_A^{\delta}\bigl( a(x) \bigr)$.
\end{lemma}

\begin{proof}
\leanok
\uses{lem:exists-point}
The average of the lower verifier exceeds $1/\delta - \Phi(A)$, which by assumption is
at least $1/\zeta + \Psi_J(B)$, which exceeds the average of the upper verifier. Now
apply \cref{lem:exists-point}.
\end{proof}

Any weight $w$ between the two reciprocals is then admissible for both halves of
\cref{lem:barrier}. The construction always takes the largest one, $w = 1/L_A^{\delta}$.

\chapter{The construction and the frame bounds}
\label{ch:construction}

\begin{definition}[The two states]
\label{def:states}
\lean{Discretization.lowerState, Discretization.upperState}
\leanok
For initial matrices $A_0$, $B_0$, shifts $\delta, \zeta$, points
$x_1, \dots, x_k$ and weights $w_1, \dots, w_k$, put
\[
  A_k = A_0 - k \delta \cdot 1 + \sum_{i=1}^k w_i \, a(x_i) a(x_i)^*,
  \qquad
  B_k = B_0 + k \zeta \cdot J - \sum_{i=1}^k w_i \, b(x_i) b(x_i)^* .
\]
\end{definition}

\begin{proposition}[The construction]
\label{prop:construction}
\lean{Discretization.exists_points_weights}
\leanok
\uses{def:states}
Let $A_0$, $B_0$ and $J$ be positive definite, let $\delta, \zeta > 0$, let $G(a) = 1$
and $G(b) = J$, and assume the initial gap condition
\[
  \frac{1}{\zeta} + \Psi_J(B_0) \; \leq \; \frac{1}{\delta} - \Phi(A_0).
\]
Then for every $k$ there are points $x_1, \dots, x_k \in \Omega$ and weights
$w_1, \dots, w_k > 0$ such that $A_k$ and $B_k$ are positive definite and
\[
  \Phi(A_k) \leq \Phi(A_0), \qquad \Psi_J(B_k) \leq \Psi_J(B_0).
\]
\end{proposition}

\begin{proof}
\leanok
\uses{lem:admissible-point, lem:barrier}
Induction on $k$. The invariant keeps the gap open, so \cref{lem:admissible-point}
produces a point $x_{k+1}$; the weight is the reciprocal of its lower verifier value,
and \cref{lem:barrier} restores the invariant.
\end{proof}

The number of steps is arbitrary in \cref{prop:construction}: only the initial gap
condition drives it. It is the choice of the four parameters that ties the number of
points to the frame bounds. With $m$ the size of the first family, $M$ the effective
dimension of the second one and $n$ the number of points to be produced, put
\[
  r = \sqrt{\tfrac{m-1}{n}}, \quad s = \sqrt{\tfrac{M-1}{n}}, \quad
  \delta = \tfrac{1-r}{n}, \quad \zeta = \tfrac{1+s}{n},
\]
\[
  A_0 = \tfrac{\delta m}{r} \cdot 1, \qquad
  B_0 = \tfrac{\zeta \Tr J}{s} \cdot 1 .
\]
These are chosen so that the initial gap closes exactly,
$1/\delta - \Phi(A_0) = n = 1/\zeta + \Psi_J(B_0)$, which needs $m \geq 2$ for $r > 0$
and $M \geq 1 + 1/n$ for $s \geq 1/n$.

\begin{lemma}[Reading off the frame bounds]
\label{lem:read-off}
\lean{Discretization.lower_bound_of_state, Discretization.upper_bound_of_state}
\leanok
\uses{def:states}
Let $A_k$ be positive definite with $\Phi(A_k) \leq c$ and $A_0 = c_0 \cdot 1$. Then
\[
  \Bigl( c^{-1} - c_0 + k \delta \Bigr) \cdot 1
    \; \preceq \; \sum_{i=1}^k w_i \, a(x_i) a(x_i)^* .
\]
Let $B_k$ be positive definite with $\Psi_J(B_k) \leq c$ and $B_0 = d_0 \cdot 1$. Then
\[
  \sum_{i=1}^k w_i \, b(x_i) b(x_i)^*
    \; \preceq \; d_0 \cdot 1 + \bigl( k \zeta - c^{-1} \bigr) \cdot J .
\]
\end{lemma}

\begin{proof}
\leanok
\uses{lem:inv-trace-smul-one-le, lem:inv-trace-mul-smul-le}
A bound on a potential is a bound on the matrix by
\cref{lem:inv-trace-smul-one-le} and \cref{lem:inv-trace-mul-smul-le}. Substituting the
definition of the state and rearranging gives the two inequalities.
\end{proof}

\begin{theorem}[Generalized sparsification, normalized first family]
\label{thm:main-normalized}
\lean{Discretization.bss_generalized_of_gram_eq_one',
      Discretization.bss_generalized_of_gram_eq_one}
\leanok
\uses{def:effective-dimension}
Let $a$ and $b$ be square-integrable in every coordinate with $G(a) = 1$ and
$G(b) = J$ positive definite, and let $\Lambda > 0$ with $J \preceq \Lambda \cdot 1$.
Put $m = \#\iota$ and $M = \Tr J / \Lambda$. Then for every $n \geq m$ there are points
$x_1, \dots, x_n \in \Omega$ and weights $w_1, \dots, w_n > 0$ with
\[
  \Bigl(1 - \sqrt{\tfrac{m-1}{n}}\Bigr)^2 \cdot 1
    \; \preceq \; \sum_{i=1}^n w_i \, a(x_i) a(x_i)^*
\]
and
\[
  \sum_{i=1}^n w_i \, b(x_i) b(x_i)^*
    \; \preceq \; \Bigl(1 + \sqrt{\tfrac{M-1}{n}}\Bigr)^2 \Lambda \cdot 1 .
\]
\end{theorem}

\begin{proof}
\leanok
\uses{prop:construction, lem:read-off, lem:frame-constants, thm:card-one, thm:small-dim,
      thm:both-edge-cases}
For $m \geq 2$ and $M \geq 1 + 1/n$, the parameters above satisfy the gap condition of
\cref{prop:construction} with equality, so $n$ steps produce points and weights with
$\Phi(A_n) \leq r/\delta$ and $\Psi_J(B_n) \leq s/\zeta$. Feeding these into
\cref{lem:read-off} and inserting the parameters through
\cref{lem:frame-constants} gives the two bounds. The cases $m = 1$ and
$M \leq 1 + 1/n$ are \cref{thm:card-one}, \cref{thm:small-dim} and
\cref{thm:both-edge-cases}, and a case distinction on the two conditions assembles the
four into the statement above.
\end{proof}

\begin{lemma}[The frame constants]
\label{lem:frame-constants}
\lean{Discretization.lower_frame_bound, Discretization.upper_frame_bound}
\leanok
With the parameters above, $c_0 = \delta m / r$, $d_0 = \zeta \Tr J / s$ and
$\Tr J = \Lambda (n s^2 + 1)$,
\[
  \frac{\delta}{r} + n \delta - c_0 = (1 - r)^2,
  \qquad
  d_0 + \Bigl( n \zeta - \frac{\zeta}{s} \Bigr) \Lambda = (1 + s)^2 \Lambda ,
\]
and the coefficient $n \zeta - \zeta/s$ of $J$ is nonnegative.
\end{lemma}

\begin{proof}
\leanok
Arithmetic with the four parameters; the square roots are never unfolded. The
nonnegativity of the coefficient is where $M \geq 1 + 1/n$, in the form $s \geq 1/n$,
is used, and it is what allows $J \preceq \Lambda \cdot 1$ to be applied to it.
\end{proof}

\chapter{The two edge cases}
\label{ch:edge}

The parameters of \cref{ch:construction} need $m \geq 2$ and $M \geq 1 + 1/n$. Where one
of them fails, the corresponding potential is not needed at all: the verifier on that
side is replaced by a constant, whose average is again $n$.

\begin{theorem}[A one-element first family]
\label{thm:card-one}
\lean{Discretization.bss_generalized_of_unique}
\leanok
Let $\#\iota = 1$, let $a$ and $b$ be square-integrable in every coordinate with
$G(a) = 1$ and $G(b) = J$ positive definite, let $\Lambda > 0$ with
$J \preceq \Lambda \cdot 1$ and let $n \geq 1$ with $M = \Tr J / \Lambda \geq 1 + 1/n$.
Then there are $n$ points and weights $w_i > 0$ with
\[
  1 \; \preceq \; \sum_{i=1}^n w_i \, a(x_i) a(x_i)^*
  \qquad\text{and}\qquad
  \sum_{i=1}^n w_i \, b(x_i) b(x_i)^*
    \; \preceq \; \Bigl(1 + \sqrt{\tfrac{M-1}{n}}\Bigr)^2 \Lambda \cdot 1 .
\]
\end{theorem}

\begin{proof}
\leanok
\uses{lem:barrier, lem:average-upper, lem:exists-point, lem:read-off}
The lower verifier is replaced by $L(x) = n |a(x)|^2$, whose average is $n$ because the
single function is normalized. The induction then tracks only the upper matrix, using
the upper half of \cref{lem:barrier}, and takes the largest weight the constant allows,
$1/w_i = L(x_i)$. Each point therefore contributes exactly $1/n$ to the lower bound, which
becomes the identity $\sum_i w_i |a(x_i)|^2 = 1 = (1 - \sqrt{(m-1)/n})^2$.
\end{proof}

\begin{theorem}[A small effective dimension]
\label{thm:small-dim}
\lean{Discretization.bss_generalized_of_small_dim}
\leanok
Let $a$ and $b$ be square-integrable in every coordinate with $G(a) = 1$ and
$G(b) = J$ positive definite, let $\Lambda > 0$, let $m = \#\iota \geq 2$ and
$n \geq m$, and assume $M = \Tr J / \Lambda \leq 1 + 1/n$. Then there are $n$ points and
weights $w_i > 0$ with
\[
  \Bigl(1 - \sqrt{\tfrac{m-1}{n}}\Bigr)^2 \cdot 1
    \; \preceq \; \sum_{i=1}^n w_i \, a(x_i) a(x_i)^*
\]
and
\[
  \sum_{i=1}^n w_i \, b(x_i) b(x_i)^*
    \; \preceq \; \Bigl(1 + \sqrt{\tfrac{M-1}{n}}\Bigr)^2 \Lambda \cdot 1 .
\]
\end{theorem}

\begin{proof}
\leanok
\uses{lem:barrier, lem:average-lower, lem:exists-point, lem:read-off, lem:rank-one-crude}
The upper verifier is replaced by $U(x) = n \|b(x)\|^2 / \Tr J$, whose average is $n$.
The induction tracks only the lower matrix, and the weights it produces satisfy
$w_i U(x_i) \leq 1$, that is $w_i \|b(x_i)\|^2 \leq \Tr J / n$. Summing over the $n$
points and using \cref{lem:rank-one-crude} gives
$\sum_i w_i \, b(x_i) b(x_i)^* \preceq \Tr J \cdot 1 = M \Lambda \cdot 1$, and
$M \leq 1 + 1/n$ forces $s \leq 1/n$ and hence $M = 1 + n s^2 \leq 1 + s \leq (1+s)^2$.
\end{proof}

\begin{theorem}[Both edge cases at once]
\label{thm:both-edge-cases}
\lean{Discretization.bss_generalized_of_unique_of_small_dim}
\leanok
Let $\#\iota = 1$, let $a$ and $b$ be square-integrable in every coordinate with
$G(a) = 1$ and $G(b) = J$ positive definite, let $\Lambda > 0$ and $n \geq 1$, and
assume $M = \Tr J / \Lambda \leq 1 + 1/n$. Then there are $n$ points and weights
$w_i > 0$ with
\[
  1 \; \preceq \; \sum_{i=1}^n w_i \, a(x_i) a(x_i)^*
  \qquad\text{and}\qquad
  \sum_{i=1}^n w_i \, b(x_i) b(x_i)^*
    \; \preceq \; \Bigl(1 + \sqrt{\tfrac{M-1}{n}}\Bigr)^2 \Lambda \cdot 1 .
\]
\end{theorem}

\begin{proof}
\leanok
\uses{lem:exists-point-le, lem:rank-one-crude}
Both verifiers are now the constants $L(x) = n |a(x)|^2$ and
$U(x) = n \|b(x)\|^2 / \Tr J$, of equal average $n$, and neither depends on the state of
the construction. \cref{lem:exists-point-le} produces a single point $y$ with
$U(y) \leq L(y)$ and $L(y) > 0$; using it $n$ times with the weight $w = 1/L(y)$ makes
the lower bound the identity $\sum_i w_i |a(x_i)|^2 = 1$, while
$w \, U(y) \leq 1$ and \cref{lem:rank-one-crude} give the upper bound as in
\cref{thm:small-dim}. No induction is needed.
\end{proof}

\chapter{A countably infinite second family}
\label{ch:infinite}

Nothing in \cref{thm:main-normalized} ties the number of points to the size of the
second family: it is governed by the effective dimension alone. The second family may
therefore be infinite. In this chapter it is a single square-integrable map
$b : \Omega \to H$ into a separable Hilbert space, and the Gram matrix becomes the
positive operator $J$ on $H$ determined by
\[
  \langle u, J u \rangle = \int_\Omega \bigl| \langle u, b(x) \rangle \bigr|^2 d\mu(x),
  \qquad u \in H .
\]
The first family stays finite, so the lower half of the argument is unchanged, and only
the upper half has to be redone. Mathlib has no trace-class theory, so the trace is
built along a fixed Hilbert basis.

\section{The trace of an operator along a basis}

\begin{definition}[Trace along a Hilbert basis]
\label{def:trace-along}
\lean{ContinuousLinearMap.traceAlong}
\leanok
Let $e = (e_k)_{k \in K}$ be a Hilbert basis of $H$ and $T$ a bounded operator on $H$.
Put
\[
  \Tr_e T = \sum_{k \in K} \re \langle e_k, T e_k \rangle,
\]
with the convention that the sum is zero when the family is not summable.
\end{definition}

No basis independence is proved and none is needed: every statement below fixes one
basis, and in the application the basis is the singular basis of the embedding. For a
positive $T$ the terms are nonnegative, because
$\re \langle x, T x \rangle = \|T^{1/2} x\|^2$, so $\Tr_e T$ is the squared
Hilbert--Schmidt norm of $T^{1/2}$ and every statement about its value assumes the
family summable.

\begin{lemma}[Parseval's identity along a basis]
\label{lem:parseval}
\lean{HilbertBasis.hasSum_norm_sq_inner}
\leanok
For every Hilbert basis $e$ of $H$ and every $x \in H$, the family
$\bigl( |\langle e_k, x \rangle|^2 \bigr)_k$ has sum $\|x\|^2$.
\end{lemma}

\begin{proof}
\leanok
Take real parts in Mathlib's Parseval identity
$\sum_k \langle x, e_k \rangle \langle e_k, x \rangle = \langle x, x \rangle$.
\end{proof}

\begin{lemma}[Invariance under adjoints]
\label{lem:hs-adjoint}
\lean{ContinuousLinearMap.tsum_norm_sq_adjoint,
      ContinuousLinearMap.summable_norm_sq_adjoint_iff}
\leanok
\uses{def:trace-along}
For every bounded operator $T$ on $H$ and every Hilbert basis $e$,
\[
  \sum_k \| T e_k \|^2 = \sum_k \| T^* e_k \|^2 ,
\]
and the family on the left is summable if and only if the one on the right is.
\end{lemma}

\begin{proof}
\leanok
\uses{lem:parseval}
Expanding both sides by \cref{lem:parseval} gives the same double series
$\sum_{j,k} |\langle e_j, T e_k \rangle|^2$ with the two indices exchanged. The
rearrangement is carried out in $[0, \infty]$, where every sum is defined and the
interchange needs no hypothesis; summability is then read off the value being finite.
\end{proof}

This is the hinge of the whole trace layer: it is what supplies the summability of
$P^{1/2} Q$, whose adjoint $Q P^{1/2}$ is visibly Hilbert--Schmidt.

\begin{lemma}[Crude bound for operators]
\label{lem:le-trace-along-smul-one}
\lean{ContinuousLinearMap.le_traceAlong_smul_one}
\leanok
\uses{def:trace-along}
Let $T$ be a positive operator on $H$ whose trace along $e$ is summable. Then
$T \preceq \Tr_e T \cdot 1$.
\end{lemma}

\begin{proof}
\leanok
\uses{lem:parseval}
For $x \in H$, writing $x$ in the basis and applying the Cauchy--Schwarz inequality to
each term gives $\re \langle x, T x \rangle \leq \Tr_e T \cdot \|x\|^2$. This is the
operator analogue of \cref{lem:le-trace-smul-one}, and unlike the matrix proof it uses
no diagonalisation, which is also why no basis independence is required.
\end{proof}

\begin{lemma}[Trace of a product]
\label{lem:trace-along-mul}
\lean{ContinuousLinearMap.traceAlong_mul, ContinuousLinearMap.traceAlong_mul_pos}
\leanok
\uses{def:trace-along}
Let $P$ be positive with summable trace along $e$ and let $Q$ be positive. Then
\[
  \Tr_e (P Q) = \sum_k \re \bigl\langle P^{1/2} e_k, \, Q \, P^{1/2} e_k \bigr\rangle
  \; \geq \; 0 ,
\]
and the value is positive as soon as $Q$ is invertible and $P \neq 0$.
\end{lemma}

\begin{proof}
\leanok
\uses{lem:hs-adjoint}
The identity is the cyclicity $\Tr (PQ) = \Tr (P^{1/2} Q P^{1/2})$, read as a symmetry
of the pairing $\sum_k \langle S e_k, T e_k \rangle$ under passing to adjoints; the
required summability is \cref{lem:hs-adjoint}. Every term on the right is nonnegative.
For the strict inequality, invertibility of $Q$ makes $Q^{1/2}$ injective, so a
vanishing trace would force $P^{1/2}$ to annihilate the whole basis.
\end{proof}

Invertibility of $Q$ here is what replaces positive definiteness. A positive operator of
finite trace on an infinite-dimensional space is compact, hence never bounded away from
zero, so positive definiteness is unavailable and the hypotheses on the two operators
have to be chosen differently.

\begin{definition}[Positive, injective, of finite trace]
\label{def:finite-trace-pos}
\lean{Discretization.Infinite.IsFiniteTracePos}
\leanok
\uses{def:trace-along}
An operator $J$ on $H$ is \emph{of finite trace and positive} relative to a Hilbert
basis $e$ if $J$ is positive, its trace along $e$ is summable, and $J$ is injective.
\end{definition}

For the second operator, the one the construction moves, the right notion is positivity
together with invertibility; it is Mathlib's \texttt{IsStrictlyPositive}, and the inverse
is \texttt{Ring.inverse}, defined on all operators and equal to the inverse on the
invertible ones.

\begin{lemma}[Sherman--Morrison for operators]
\label{lem:sherman-morrison-op}
\lean{ContinuousLinearMap.inverse_add_smul_rankOne}
\leanok
Let $A$ be a self-adjoint invertible operator on $H$, let $u \in H$ and let $t$ be a
scalar with $1 + t \langle u, A^{-1} u \rangle \neq 0$. Then $A + t \, u u^*$ is
invertible and
\[
  \bigl( A + t \, u u^* \bigr)^{-1}
    = A^{-1} - \frac{t}{1 + t \langle u, A^{-1} u \rangle} \,
      \bigl( A^{-1} u \bigr) \bigl( A^{-1} u \bigr)^* ,
\]
where $u u^*$ is the rank-one operator $z \mapsto \langle u, z \rangle u$.
\end{lemma}

\begin{proof}
\leanok
The candidate is verified to be a two-sided inverse, exactly as in
\cref{lem:sherman-morrison}. In a noncommutative ring one product does not suffice, but
two do: they exhibit $A + t \, u u^*$ as a unit.
\end{proof}

\section{The upper potential of an operator}

\begin{definition}[The upper potential of an operator]
\label{def:upper-potential-op}
\lean{Discretization.Infinite.upperPotential}
\leanok
\uses{def:trace-along, def:finite-trace-pos}
For $J$ of finite trace and positive, and $B$ strictly positive,
$\Psi_J(B) = \Tr_e \bigl( J B^{-1} \bigr)$.
\end{definition}

\begin{lemma}[Effect of the shift]
\label{lem:shift-effect-op}
\lean{Discretization.Infinite.upperPotential_sub_eq,
      Discretization.Infinite.upperPotential_add_smul_lt}
\leanok
\uses{def:upper-potential-op}
Let $J$ be of finite trace and positive, let $B$ be strictly positive and
$\zeta \geq 0$. Then
\[
  \Psi_J(B) - \Psi_J(B + \zeta \cdot J)
    = \zeta \, \Tr_e \bigl( (J B^{-1} J) (B + \zeta \cdot J)^{-1} \bigr),
\]
and for $\zeta > 0$ the right-hand side is positive.
\end{lemma}

\begin{proof}
\leanok
\uses{lem:trace-along-mul}
The resolvent identity holds in any ring. The right-hand side is the trace of a product
of two positive operators, the first of them $J B^{-1} J$, which is again of finite
trace and positive; \cref{lem:trace-along-mul} gives the sign, and injectivity of $J$
makes it strict.
\end{proof}

\begin{lemma}[Average of an operator quadratic form]
\label{lem:integral-inner-op}
\lean{ContinuousLinearMap.integral_re_inner_apply}
\leanok
\uses{def:trace-along}
Let $b : \Omega \to H$ be square-integrable with Gram operator $J$ of finite trace, and
let $Q$ be a positive bounded operator. Then
\[
  \int_\Omega \re \bigl\langle b(x), Q \, b(x) \bigr\rangle \, d\mu(x)
    = \Tr_e (J Q).
\]
\end{lemma}

\begin{proof}
\leanok
\uses{lem:parseval, lem:hs-adjoint, lem:trace-along-mul}
Pointwise $\re \langle b, Q b \rangle = \|Q^{1/2} b\|^2$, which \cref{lem:parseval}
turns into $\sum_k |\langle Q^{1/2} e_k, b \rangle|^2$. Sum and integral may be
exchanged because all terms are nonnegative and their integrals sum to a finite value,
and the Gram identity turns the result into
$\sum_k \re \langle Q^{1/2} e_k, J \, Q^{1/2} e_k \rangle$. That equals $\Tr_e (J Q)$
because both are the squared Hilbert--Schmidt norm of $Q^{1/2} J^{1/2}$, once read
through the operator and once through its adjoint, by \cref{lem:hs-adjoint} and
\cref{lem:trace-along-mul}. Countability of the index set of the basis enters here, in
the interchange.
\end{proof}

\section{The barrier lemma, the bound and the theorem}

\begin{lemma}[Barrier lemma for operators]
\label{lem:barrier-op}
\lean{Discretization.Infinite.upperVerifier,
      Discretization.Infinite.upperPotential_update_le}
\leanok
\uses{def:upper-potential-op}
Write $X = (B + \zeta \cdot J)^{-1}$ and
\[
  U_B^{\zeta}(u)
    = \frac{\re \langle u, X J X u \rangle}
           {\Psi_J(B) - \Psi_J(B + \zeta \cdot J)} + \re \langle u, X u \rangle .
\]
Let $J$ be of finite trace and positive, let $B$ be strictly positive, $\zeta > 0$ and
let $w > 0$ satisfy $U_B^{\zeta}(u) \leq 1/w$. Then
$B + \zeta \cdot J - w \, u u^*$ is strictly positive and
$\Psi_J \bigl( B + \zeta \cdot J - w \, u u^* \bigr) \leq \Psi_J(B)$.
\end{lemma}

\begin{proof}
\leanok
\uses{lem:sherman-morrison-op, lem:shift-effect-op, lem:barrier}
The inequality between real numbers is the one of \cref{lem:barrier}; what has to be
redone is the closed form of the potential after a rank-one downdate, which rests on
\cref{lem:sherman-morrison-op}, and the positivity of the numerator, which rests on the
injectivity of $J$.
\end{proof}

\begin{lemma}[The verifier passes on average]
\label{lem:average-upper-op}
\lean{Discretization.Infinite.integral_upperVerifier_lt,
      Discretization.Infinite.exists_admissible_point}
\leanok
\uses{lem:barrier-op}
In the situation of \cref{lem:barrier-op}, with $b$ square-integrable and $J$ its Gram
operator,
\[
  \int_\Omega U_B^{\zeta} \bigl( b(x) \bigr) d\mu(x) < \frac{1}{\zeta} + \Psi_J(B).
\]
Consequently, if the gap condition
$1/\zeta + \Psi_J(B) \leq 1/\delta - \Phi(A)$ holds for a positive definite matrix $A$
and the first family, then some point passes the test of both verifiers.
\end{lemma}

\begin{proof}
\leanok
\uses{lem:integral-inner-op, lem:average-lower, lem:exists-point}
\cref{lem:integral-inner-op} turns the average into
$\Tr_e (J X J X) / E + \Psi_J(B + \zeta \cdot J)$ with $E$ the gain of
\cref{lem:shift-effect-op}, and $X \preceq B^{-1}$ bounds the numerator by $E/\zeta$.
The comparison with the lower verifier of \cref{lem:average-lower} is a comparison of
two averages of real functions, one coming from a matrix and one from an operator, so
\cref{lem:exists-point} applies unchanged.
\end{proof}

\begin{lemma}[A bound on the potential bounds the operator]
\label{lem:inv-potential-op}
\lean{Discretization.Infinite.inv_upperPotential_smul_le}
\leanok
\uses{def:upper-potential-op}
Let $J$ be of finite trace and positive and let $B$ be strictly positive. Then
$\Psi_J(B)^{-1} \cdot J \preceq B$.
\end{lemma}

\begin{proof}
\leanok
\uses{lem:le-trace-along-smul-one, lem:sqrt-conjugation}
The operator $B^{-1/2} J B^{-1/2}$ is positive and its trace along $e$ is again
$\Psi_J(B)$, so \cref{lem:le-trace-along-smul-one} bounds it by $\Psi_J(B) \cdot 1$, and
\cref{lem:sqrt-conjugation} carries that back to $J \preceq \Psi_J(B) \cdot B$.
\end{proof}

\begin{theorem}[Sparsification with a countable second family]
\label{thm:main-infinite}
\lean{Discretization.Infinite.bss_generalized_of_gram_eq_one}
\leanok
\uses{def:finite-trace-pos}
Let $a : \Omega \to \mathbb{C}^{\iota}$ be square-integrable in every coordinate with
$G(a) = 1$ and $m = \#\iota \geq 2$, and let $b : \Omega \to H$ be square-integrable
with Gram operator $J$ of finite trace, positive and injective, and
$J \preceq \Lambda \cdot 1$ for some $\Lambda > 0$. Put $M = \Tr_e J / \Lambda$ and
assume $M \geq 1 + 1/n$. Then for every $n \geq m$ there are points
$x_1, \dots, x_n \in \Omega$ and weights $w_1, \dots, w_n > 0$ with
\[
  \Bigl(1 - \sqrt{\tfrac{m-1}{n}}\Bigr)^2 \cdot 1
    \; \preceq \; \sum_{i=1}^n w_i \, a(x_i) a(x_i)^*
\]
in the matrix order and
\[
  \sum_{i=1}^n w_i \, b(x_i) b(x_i)^*
    \; \preceq \; \Bigl(1 + \sqrt{\tfrac{M-1}{n}}\Bigr)^2 \Lambda \cdot 1
\]
in the order of operators.
\end{theorem}

\begin{proof}
\leanok
\uses{lem:barrier-op, lem:average-upper-op, lem:inv-potential-op, prop:construction,
      lem:frame-constants, lem:read-off}
The construction is that of \cref{prop:construction} with the upper state the operator
$B_k = B_0 + k \zeta \cdot J - \sum_i w_i \, b(x_i) b(x_i)^*$: the step is
\cref{lem:average-upper-op} followed by the two barrier lemmas, the lower one for
matrices and \cref{lem:barrier-op} for operators. The initial data, the parameters and
the lower read-off are those of \cref{ch:construction} verbatim; the upper read-off is
\cref{lem:inv-potential-op} in place of \cref{lem:inv-trace-mul-smul-le}, and the
constants come out the same by \cref{lem:frame-constants}.
\end{proof}

\begin{corollary}[Discretization with a countable second family]
\label{cor:discretization-infinite}
\lean{Discretization.Infinite.exists_discretization}
\leanok
\uses{thm:main-infinite}
In the situation of \cref{thm:main-infinite}, the same points and weights satisfy, for
every $c \in \mathbb{C}^{\iota}$,
\[
  \Bigl(1 - \sqrt{\tfrac{m-1}{n}}\Bigr)^2 \int_\Omega
    \Bigl| \sum_k \overline{c_k} \, a_k(x) \Bigr|^2 d\mu(x)
  \; \leq \; \sum_{i=1}^n w_i \Bigl| \sum_k \overline{c_k} \, a_k(x_i) \Bigr|^2
\]
and, for every $u \in H$,
\[
  \sum_{i=1}^n w_i \bigl| \langle u, b(x_i) \rangle \bigr|^2
  \; \leq \; \Bigl(1 + \sqrt{\tfrac{M-1}{n}}\Bigr)^2 \Lambda \, \|u\|^2 .
\]
\end{corollary}

\begin{proof}
\leanok
\uses{lem:integral-quadform}
A Loewner inequality between operators is an inequality between quadratic forms,
uniformly in the vector; for the first family this is \cref{lem:integral-quadform} read
backwards, and for the second the quadratic form of a weighted sum of rank-one operators
is the weighted sum of the squared moduli.
\end{proof}

\begin{remark}[What is left]
\label{rem:left}
\leanok
The number of points in \cref{thm:main-infinite} is governed by the effective dimension
$M = \Tr_e J / \Lambda$ and not by the size of the second family. What is not formalised
is the construction of the Gram operator from a reproducing kernel: the singular value
decomposition of the embedding $H \hookrightarrow L_2(\mu)$, which supplies the family
$b$ and makes $J$ diagonal, with $\Tr_e J$ the trace of the kernel and $\Lambda$ the
squared norm of the embedding.
\end{remark}

\bpbibheading
\begin{thebibliography}{9}

\bibitem{BSS}
J.~Batson, D.~A.~Spielman, N.~Srivastava,
\emph{Twice-Ramanujan sparsifiers},
SIAM Review \textbf{56} (2014), 315--334.

\bibitem{CDKU}
A.~Chkifa, M.~Dolbeault, D.~Krieg, M.~Ullrich,
\emph{Constructive discretization and approximation in reproducing kernel Hilbert
spaces}.

\end{thebibliography}
